\section*{Capítulo \ref{ch:b2:sound-waves} --- Ondas sonoras en los fluidos}

\begin{solution}{exo:b2:sound-waves:1}
$c = \sqrt{\gamma RT/M}$: \qty{331}{m/s}, \qty{343}{m/s}, \qty{299}{m/s}; helio
\qty{1010}{m/s}; CO$_2$ \qty{268}{m/s}. Las cuerdas vocales vibran a la misma
frecuencia, pero las resonancias del tracto vocal (que dan forma al
timbre) escalan con $c$ y suben en un factor tres: la voz suena aguda.
\end{solution}

\begin{solution}{exo:b2:sound-waves:2}
\qty{60}{dB}: $I = \qty{1e-6}{W/m^2}$, $p_m = \sqrt{2ZI} = \qty{29}{mPa}$, $v_m = p_m/Z =
\qty{7e-5}{m/s}$, $\xi_m = v_m/\omega = \qty{22}{nm}$. \qty{120}{dB}: \qty{1}{W/m^2},
\qty{29}{Pa}, \qty{7}{cm/s}, \qty{22}{\micro m}. Dos: \qty{63}{dB}; cien: \qty{80}{dB}.
\end{solution}

\begin{solution}{exo:b2:sound-waves:3}
$I = 1/4\pi r^2$: \qty{0.080}{W/m^2} (\qty{109}{dB}), \qty{89}{dB}, \qty{69}{dB}; $I = 10^{-6}$
en $r = \sqrt{1/4\pi10^{-6}} = \qty{280}{m}$; para \qty{100}{dB} a \qty{30}{m} hacen falta $4\pi \times
900 \times 10^{-2} = \qty{110}{W}$ de sonido.
\end{solution}

\begin{solution}{exo:b2:sound-waves:4}
$4.5 \times 340 = \qty{1.5}{km}$; $340 \times 1.5 = \qty{510}{m}$; $1500 \times 1.2 = \qty{1.8}{km}$;
$1540 \times 65 \times 10^{-6} = \qty{10}{cm}$.
\end{solution}

\begin{solution}{exo:b2:sound-waves:5}
(a) $500 \times 340/300 = \qty{567}{Hz}$; $500 \times 340/380 = \qty{447}{Hz}$. (b) $500 \times
380/340 = \qty{559}{Hz}$; ambos: $500 \times 380/300 = \qty{633}{Hz}$, frente a \qty{654}{Hz}
(solo la fuente a \qty{80}{m/s}) o \qty{618}{Hz} (solo el observador): lo que fija
el resultado es el aire, no la celeridad relativa. (c) La pared recibe \qty{567}{Hz}
y la reemite en reposo; el maquinista se acerca a \qty{40}{m/s}: $567 \times
380/340 = \qty{633}{Hz}$.
\end{solution}

\begin{solution}{exo:b2:sound-waves:6}
$v_m = p_m/Z$: \qty{2.4}{mm/s}, \qty{0.67}{\micro m/s}, \qty{25}{nm/s}; $\xi_m = v_m/\omega$:
\qty{390}{nm}, \qty{0.11}{nm}, \qty{4}{pm}; $I = p_m^2/2Z$: \qty{1.2e-3}{W/m^2} (\qty{91}{dB}),
\qty{3.3e-7}{W/m^2} (\qty{55}{dB}), \qty{1.3e-8}{W/m^2} (\qty{41}{dB}). Un medio rígido
apenas se mueve bajo una presión dada y transporta poca potencia.
\qty{1}{Pa} son \qty{120}{dB} re \qty{1}{\micro Pa}.
\end{solution}

\begin{solution}{exo:b2:sound-waves:7}
(a) $\sqrt{1.013 \times 10^5/1.2} = \qty{290}{m/s}$; $\times\sqrt{1.4}$: \qty{343}{m/s}. (b)
$\gamma = c^2M/RT$: $1.40$ para el aire, $1.71$ para el argón ($5/3$ dentro del redondeo):
tres grados de libertad de traslación ($\gamma = 1 + 2/3$) frente a cinco para
un gas diatómico ($1 + 2/5$). (c) $\sqrt{D/f} = \sqrt{2 \times 10^{-8}} = \qty{0.14}{mm}
\ll \lambda = \qty{34}{cm}$: no pasa calor entre crestas y valles.
\end{solution}

\begin{solution}{exo:b2:sound-waves:8}
(a) Abierto--abierto: $p = p_m\sin kx\cos\omega t$ con $kL = n\pi$; cerrado--abierto:
$p = p_m\cos kx\cos\omega t$ (nodo de velocidad y vientre de presión en el extremo
cerrado) con $kL = (2n - 1)\pi/2$. (b) $L = c/2f = \qty{39}{cm}$; $c/4f = \qty{19.5}{cm}$.
(c) Todos los armónicos; solo los impares (el sonido hueco del clarinete).
(d) $0.6 \times 2 = \qty{1.2}{cm}$ por extremo abierto: $L_{\text{ef}} = \qty{41.4}{cm}$, $f =
\qty{414}{Hz}$, un $6\%$ bajo --- hay que cortar el tubo a \qty{36.6}{cm}.
\end{solution}

\begin{solution}{exo:b2:sound-waves:9}
$\mathcal P = \qty{400}{W}$, $I = \mathcal P/2\pi r^2$. (a) \qty{6.4e-3}{W/m^2}, \qty{98}{dB},
$p_m = \qty{2.3}{Pa}$. (b) $r = \sqrt{400/2\pi} = \qty{8}{m}$; \qty{70}{dB} a \qty{2.5}{km}.
(c) $98 - 20\log(r/100) - (r - 100)/100 = 70$: $r \approx \qty{950}{m}$. (d)
$\pi(10)^2/\pi(0.95)^2 \approx 110$ sirenas.
\end{solution}

\begin{solution}{exo:b2:sound-waves:10}
(a) $\rho_0\partial_tv = p_mk\sin kx\cos\omega t$: $v = (p_m/Z)\sin kx\sin\omega t$. (b)
$e_k = \tfrac12\chi_Sp_m^2\sin^2kx\sin^2\omega t$, $e_p = \tfrac12\chi_Sp_m^2\cos^2kx\cos^2\omega t$:
una es máxima cuando la otra se anula, en el tiempo y en el espacio. (c) $E =
S\int_0^L(e_k + e_p)\dd x = \tfrac14\chi_Sp_m^2SL$, constante. (d) $I = pv =
(p_m^2/4Z)\sin2kx\sin2\omega t$: de media nula, pero dos veces por periodo la
energía va de los vientres de presión a los de velocidad y vuelve.
\end{solution}

\begin{solution}{exo:b2:sound-waves:11}
(a) El frente emitido en $t = 0$ tiene radio $ct$ cuando la fuente está a $vt$
de distancia: el cono tangente cumple $\sin\theta = c/v$. (b) $\theta = \qty{39}{\degree}$; el cono
llega al suelo bajo el avión cuando este se ha alejado $h/\tan\theta = \qty{15}{km}$
más: $15000/480 = \qty{31}{s}$ después del paso por la vertical. (c) Una franja de
\qty{60}{km} de ancho a lo largo de toda la ruta oye el estampido. (d) $\sin\theta
= 340/800$, $\theta = \qty{25}{\degree}$; primero llega el chasquido del cono y
después la detonación de la boca, que viaja a $c$ desde el arma.
\end{solution}

\begin{solution}{exo:b2:sound-waves:12}
(a) $f_1 = f(c + v\cos\theta)/c$. (b) $f_2 = f_1c/(c - v\cos\theta) = f(c + v\cos\theta)
/(c - v\cos\theta) \approx f(1 + 2v\cos\theta/c)$. (c) $2 \times 5 \times 10^6 \times 0.5 \times 0.5/1540
= \qty{1.6}{kHz}$; nula a \qty{90}{\degree}: hay que inclinar la sonda. (d) $\Delta v = 50 \times 1540/
(2 \times 5 \times 10^6 \times 0.5) = \qty{1.5}{cm/s}$. El oído sigue el tono y el timbre en
tiempo real: un silbido limpio indica \omterm{def:b2:viscous-flows:reynolds}{flujo laminar} y un siseo, turbulencia
tras un estrechamiento.
\end{solution}

\begin{solution}{pb:b2:sound-waves:1}
\textbf{1.} $c = \sqrt{1.4 \times 8.31 \times 288/0.029} = \qty{340}{m/s}$; \qty{349}{m/s}
a \qty{30}{\degreeCelsius}. \qty{1}{km} en \qty{2.9}{s}.

\textbf{2.} \qty{1.4}{km}. El canal del rayo mide kilómetros; el sonido de
sus distintas partes llega a lo largo de segundos, con ecos en las nubes y
en el suelo.

\textbf{3.} $I = \qty{1}{W/m^2}$: $p_m = \sqrt{2ZI} = \qty{29}{Pa}$, $v_m = \qty{7}{cm/s}$.

\textbf{4.} $It = \qty{2}{J/m^2}$.

\textbf{5.} $c$ crece hacia arriba; con $\sin i/c$ constante, $i$ crece, de
modo que el rayo se curva hacia la horizontal y vuelve a bajar: el sonido
queda canalizado junto al suelo y llega lejos. Sobre el agua, el aire más
bajo lo enfría el agua: lo mismo.

\textbf{6.} Los rayos se curvan hacia arriba y dejan una zona de sombra.
Radio de curvatura $R = c/(\dd c/\dd z) = \qty{50}{km}$: tras \qty{3}{km} el
rayo ha girado $3/50 = \qty{0.06}{rad} = \qty{3.4}{\degree}$ y ha subido
\qty{90}{m}: el oyente queda por debajo de él.

\textbf{7.} La absorción crece como $f^2$: las frecuencias altas del
chasquido desaparecen tras unos kilómetros y queda el retumbo grave.

\textbf{8.} $\mathcal P = \qty{2}{W}$; $I = 2/2\pi r^2 = \qty{3.2e-3}{W/m^2}$ a \qty{10}{m}:
\qty{95}{dB}.

\textbf{9.} $p_m = \sqrt{2 \times 410 \times 3.2 \times 10^{-3}} = \qty{1.6}{Pa}$; $v_m = \qty{3.9}{mm/s}$;
$\xi_m = v_m/2\pi f = \qty{6}{\micro m}$.

\textbf{10.} $+\qty{20}{dB}$: \qty{115}{dB} a \qty{1}{m} --- dañino en unos
minutos. \qty{60}{dB}: $r = \sqrt{2/2\pi10^{-6}} = \qty{560}{m}$.

\textbf{11.} Diez incoherentes: $+\qty{10}{dB}$, \qty{105}{dB}. Dos en fase:
la presión se duplica, $+\qty{6}{dB}$, \qty{101}{dB}.

\textbf{12.} Balance $\mathcal P = \alpha A\,ce/4$: $e = 8/(0.3 \times 2000 \times 340) =
\qty{3.9e-5}{J/m^3}$, $E = eV = \qty{0.4}{J}$. Tras la parada, $\dd E/\dd t = -(\alpha Ac
/4V)E$: $\tau = 4V/\alpha Ac = \qty{0.20}{s}$, y los \qty{60}{dB} tardan $\tau\ln10^6 =
\qty{2.7}{s}$ (el tiempo de reverberación de Sabine).

\textbf{13.} $\lambda = c/f$: \qty{0.51}{mm}, \qty{0.31}{mm}, \qty{0.15}{mm}.

\textbf{14.} $2 \times 0.12/1540 = \qty{156}{\micro s}$; \qty{6.4}{kHz}; $200$ líneas en
\qty{31}{ms} (treinta imágenes por segundo).

\textbf{15.} Tejido--aire: $r = -0.9995$, se refleja el $99.9\%$ de la
energía --- de ahí el gel, que elimina la película de aire.
Tejido--hueso: $r = 0.63$, el $40\%$ de la energía; el resto lo absorbe el
hueso, que proyecta una sombra. Grasa--músculo: $r = 0.10$, el $1\%$ de la
energía: los ecos tenues con los que se forma la imagen.

\textbf{16.} Fracción reflejada $1.1\%$ ($-\qty{20}{dB}$); ida y vuelta de
\qty{4}{cm} a \qty{5}{MHz}: $-\qty{10}{dB}$; en total $-\qty{30}{dB}$, una milésima.

\textbf{17.} Ida y vuelta de \qty{24}{cm}: \qty{36}{dB}, \qty{60}{dB}, \qty{120}{dB}: \qty{3}{MHz}
y \qty{5}{MHz} sirven, \qty{10}{MHz} no --- los órganos profundos se
exploran a baja frecuencia con resolución basta, y los superficiales a
alta frecuencia con detalle.

\textbf{18.} Cinco longitudes de onda, \qty{1.5}{mm} de largo: dos
interfaces separadas menos de unos \qty{0.8}{mm} (medio pulso) dan ecos
solapados.

\textbf{19.} $2 \times 5 \times 10^6 \times 0.4 \times 0.5/1540 = \qty{1.3}{kHz}$: un silbido cuyo
tono sube y baja con cada latido.

\textbf{20.} $v = \qty{25}{m/s}$: $700 \times 340/315 = \qty{756}{Hz}$, $700 \times 340/365 =
\qty{652}{Hz}$; razón $1.16$, $12\log_21.16 = 2.6$ semitonos.

\textbf{21.} $700 \times 355/315 = \qty{789}{Hz}$.

\textbf{22.} En el momento de máxima proximidad, cuando la velocidad de la
ambulancia es perpendicular a la línea de visión (en rigor, cuando el
sonido que llega ahora se emitió en ese punto).

\textbf{23.} $\lambda' = (340 - 25)/700 = \qty{0.45}{m}$ (frente a \qty{0.49}{m}); el
sonido se aleja de la ambulancia a \qty{315}{m/s} por delante y a \qty{365}{m/s}
por detrás.

\textbf{24.} $v = c\,\Delta f/2f = 3 \times 10^8 \times 3200/4.8 \times 10^{10} = \qty{20}{m/s} =
\qty{72}{km/h}$.

\textbf{25.} \qty{340}{m/s}, \qty{1540}{m/s}, \qty{25}{m/s}, \qty{20}{m/s}, \qty{0.4}{m/s},
\qty{3e8}{m/s}: $f' = f(c + v_o)/(c - v_s)$, aplicada una o dos veces, y para
la luz a baja velocidad.
\end{solution}
